\documentclass[12pt]{article}

%% Page geometry
\usepackage[margin=1in]{geometry}

%% Fonts and encoding
\usepackage[T1]{fontenc}
\usepackage[utf8]{inputenc}
\usepackage{lmodern}
\usepackage{microtype}

%% Math
\usepackage{amsmath}
\usepackage{amssymb}
\usepackage{amsthm}
\usepackage{bm}

%% Tables
\usepackage{booktabs}
\usepackage{longtable}
\usepackage{multirow}
\usepackage{array}
\usepackage{dcolumn}

%% Figures
\usepackage{graphicx}
\usepackage{float}
\usepackage{subcaption}

%% Hyperlinks and references
\usepackage[hidelinks]{hyperref}
\usepackage{natbib}
\bibliographystyle{apalike}

%% Misc
\usepackage{setspace}
\usepackage{parskip}
\usepackage{xcolor}
\usepackage{enumitem}
\usepackage{footmisc}

%% Line spacing
\setstretch{1.5}

%% Custom commands
\newcommand{\ATT}{\mathrm{ATT}}
\newcommand{\E}{\mathbb{E}}
\newcommand{\R}{\mathbb{R}}
\newcommand{\indic}{\mathbf{1}}

%% Title
\title{Wildfire Risk and Residential Property Values:\\
Evidence from the 2023 Lahaina Fire}

\author{[Author TBD]\\
\small [Affiliation TBD]\\
\small \href{mailto:}{[email TBD]}}

\date{Working Paper --- \today}

%% Abstract environment
\renewenvironment{abstract}
  {\small\begin{center}\textbf{Abstract}\end{center}\quotation}
  {\endquotation}

%%=============================================================
\begin{document}
%%=============================================================

\maketitle

\begin{abstract}
We estimate the causal effect of the August 2023 Lahaina wildfire on residential
property values in Maui County, Hawaii. Combining parcel-level transaction data
with fire perimeter shapefiles, Wildland-Urban Interface (WUI) classifications,
and aggregate price indices, we apply a multi-method identification strategy:
hedonic pricing, staggered difference-in-differences \citep{CS2021},
triple-difference decomposition, synthetic control methods
\citep{Abadie2010,BenMichael2021,Xu2017}, and spatial panel models
\citep{LeSage2009}. Our triple-difference estimates isolate a
\emph{belief-update channel} --- a forward-looking risk discount on WUI-classified
parcels beyond the physical destruction zone. Spatial autoregressive and
Spatial Durbin models quantify indirect spillovers onto properties beyond the
fire perimeter, and geographically weighted regression reveals spatial
heterogeneity in the price response. Synthetic control estimates using Maui
ZIP-level price indices corroborate the parcel-level findings. Our results
contribute to the emerging literature on how climate risk is capitalized into
housing markets \citep{Baldauf2020,Bernstein2019}.

\smallskip
\noindent\textbf{Keywords:} wildfire risk, hedonic pricing, difference-in-differences,
synthetic control, spatial econometrics, climate risk, Lahaina 2023

\smallskip
\noindent\textbf{JEL Codes:} R31, Q54, C21, C23
\end{abstract}

\newpage
\tableofcontents
\newpage

%%=============================================================
\section{Introduction}
\label{sec:intro}
%%=============================================================

The August 8, 2023 Lahaina wildfire was the deadliest United States wildfire
in more than a century, destroying approximately 2,170 structures and displacing
thousands of residents. Beyond its immediate human toll, the disaster constitutes
a natural experiment in how a catastrophic climate event is capitalized into
residential property values --- both within the destruction zone and in surrounding
areas that face elevated perceived risk.

This paper asks three questions. First, what is the average treatment effect of
fire proximity on residential sale prices in Maui County? Second, through which
economic channels does this effect operate: physical destruction, market liquidity
disruption, or a forward-looking revision of wildfire risk beliefs? Third, how far
do price spillovers extend beyond the fire perimeter, and how heterogeneous are
they across space?

Answering these questions requires meeting several econometric challenges. Sales
in the treatment area are thin following the fire, which limits the precision of
parcel-level event studies. Treatment timing is effectively uniform (a single date),
so staggered DiD methods are applied at the information-diffusion level --- using
ZIP-level insurance repricing and risk-reclassification events as cohort assignment.
Finally, spatial econometric methods are needed because parcel-level prices are
spatially autocorrelated and the ``treatment'' (perceived wildfire risk) decays
with distance in a way that induces spatial dependence in the error term.

We address these challenges with a four-phase empirical strategy. Phase 1
establishes the baseline hedonic price gradient and estimates average treatment
effects using the \citet{CS2021} doubly-robust DiD estimator. A
triple-difference model then separates the belief-update channel (WUI-classified
parcels with elevated forward-looking risk) from the physical-destruction channel.
Phase 2 constructs synthetic control counterfactuals for the Maui ZIP-level House
Price Index using the Abadie-Diamond-Hainmueller method \citep{Abadie2010},
generalized synthetic control \citep{Xu2017}, and the augmented variant
\citep{BenMichael2021}, with permutation inference to assess statistical
significance. Phase 3 estimates spatial autoregressive (SAR), spatial error (SEM),
and Spatial Durbin (SDM) models on the parcel panel, recovering direct and indirect
treatment effects via the \citet{LeSage2009} decomposition. Geographically weighted
regression (GWR) further documents spatial heterogeneity in the price response.
Phase 4 extends the data infrastructure to include FHFA ZIP-level repeat-sales HPI
and Redfin neighborhood market tracker data for robustness.

The paper proceeds as follows. Section~\ref{sec:data} describes the data.
Section~\ref{sec:empirical} presents the empirical strategy.
Section~\ref{sec:spatial} describes the spatial spillover analysis.
Section~\ref{sec:results} reports results.
Section~\ref{sec:conclusion} concludes.

%%=============================================================
\section{Data}
\label{sec:data}
%%=============================================================

\subsection{Maui County Parcel Sales}

The primary microdata source is the Maui County Real Property Assessment Roll,
which records all taxable parcel transactions in Maui County including sale price,
sale date, parcel attributes (square footage, year built, zoning category, lot size,
number of bedrooms and bathrooms, waterfront indicator), and spatial identifiers
(parcel centroid latitude and longitude, census block GEOID). The data span [TBD]
and cover approximately [TBD] residential transactions.

We restrict the analysis to arm's-length residential sales (zoning categories
R-1 through R-5 and condominium) with nominal sale prices between \$50{,}000
and \$10{,}000{,}000. All prices are deflated to [TBD] dollars using the Hawaii
component of the FHFA All-Transactions House Price Index.

\subsection{2023 Lahaina Fire Perimeter}

The Lahaina fire perimeter polygon is obtained from the National Interagency
Fire Center (NIFC) Wildfire Integrated Geographic Information System (WFIGS)
ArcGIS FeatureServer (\texttt{src/ingest/fire.py}). The perimeter represents the
final burn scar boundary as of [TBD], covering approximately 2{,}170 acres.

Distance from each parcel centroid to the fire perimeter is computed as the
minimum geodesic distance to any point on the perimeter polygon boundary,
using a vectorized nearest-neighbor query against the H3-indexed perimeter
representation (H3 resolution 8, hexagon edge length $\approx$ 460 meters).
Treatment bands are defined at 0--2 km, 2--5 km, and 5--10 km, with parcels
beyond 10 km serving as the control group.

\subsection{Wildland-Urban Interface Classification}

WUI classification is obtained from the USDA Forest Service dataset
RDS-2015-0047-3 \citep[Silvis Lab, University of Wisconsin--Madison]{},
which classifies all Census block groups in the contiguous United States as
WUI (interface or intermix), non-WUI urban, or rural. We merge WUI
classification to parcels using the census block group GEOID embedded in
the assessor data.

WUI status is used as a third dimension in the triple-difference estimator
to identify forward-looking risk beliefs. WUI-classified parcels face elevated
structural exposure to future wildfire risk, so a discount on WUI parcels relative
to non-WUI parcels within the same distance band is interpreted as a belief-update
channel rather than a physical-damage channel.

\subsection{FRED Macroeconomic Controls}

The following macroeconomic series are obtained from the St.\ Louis Fed FRED API
(\texttt{src/ingest/fred.py}) and used as time-varying controls or deflators:

\begin{itemize}[noitemsep]
  \item Hawaii median household income (\texttt{MEHOINUSHAWIA672N})
  \item FHFA All-Transactions House Price Index, Hawaii (\texttt{HISTHPI})
  \item Case-Shiller National HPI (\texttt{CSUSHPINSA})
  \item U.S. unemployment rate (\texttt{UNRATE})
  \item Effective federal funds rate (\texttt{FEDFUNDS})
  \item 30-year fixed mortgage rate (\texttt{MORTGAGE30US})
\end{itemize}

\subsection{FHFA ZIP-Level House Price Index}

The FHFA publishes a quarterly repeat-sales House Price Index at the 5-digit ZIP
code level (all-transactions, purchase-only). Data are downloaded as an Excel
workbook from the FHFA Data Tools portal (\texttt{src/ingest/fred.py},
\texttt{fetch\_fhfa\_zip\_hpi()}). These series serve as the outcome variable for
the Phase 2 synthetic control analysis, where each Maui ZIP code is the treated
unit and mainland Hawaii and continental U.S.\ ZIP codes form the donor pool.

\subsection{Redfin Neighborhood Market Tracker}

Redfin Research publishes monthly neighborhood-level housing market statistics
(median sale price, median days on market, sale-to-list price ratio, active
listings, homes sold) as a public bulk download from an S3 endpoint
(\texttt{src/ingest/redfin.py}, \texttt{fetch\_redfin\_neighborhood()}).
These data provide higher-frequency price signals and inventory indicators
that complement the annual assessment roll, and are used for robustness
checks on the hedonic price series.

%%=============================================================
\section{Empirical Strategy}
\label{sec:empirical}
%%=============================================================

\subsection{Identification}

The core identification assumption is parallel trends: absent the Lahaina fire,
treated and control parcels would have exhibited parallel log price trajectories.
We evaluate this assumption via a pre-trend regression on all pre-August-2023
event periods and by plotting event-study estimates from the Callaway-Sant'Anna
estimator.

\subsection{Hedonic Model}

Following \citet{Rosen1974}, we specify the baseline log-price hedonic:

\begin{equation}
  \log P_{it} = \alpha + \beta X_{it} + \gamma_b + \tau_t + \varepsilon_{it}
  \label{eq:hedonic}
\end{equation}

where $X_{it}$ includes log square footage, parcel age, bedroom and bathroom
counts, log lot size, zoning indicators, and a waterfront dummy; $\gamma_b$ are
census-block fixed effects; $\tau_t$ are year-month fixed effects; and
$\varepsilon_{it}$ is the HC3 heteroskedasticity-robust error \citep{MacKinnon1985}.

Treatment-band indicators $D^k_{it}$ are added to equation~(\ref{eq:hedonic})
to recover average price impacts at each distance band:

\begin{equation}
  \log P_{it} = \alpha + \beta X_{it} + \gamma_b + \tau_t + \sum_{k} \delta_k D^k_{it} + \varepsilon_{it}
  \label{eq:hedonic_treatment}
\end{equation}

\subsection{Callaway-Sant'Anna DiD}

We estimate cohort-period average treatment effects following \citet{CS2021}:

\begin{equation}
  \ATT(g,t) = \E\bigl[Y_t(g) - Y_t(0) \;\big|\; G_i = g\bigr]
  \label{eq:att}
\end{equation}

where $G_i$ is the cohort (first-treated period) of unit $i$. Estimation uses
the doubly-robust IPW estimator with a not-yet-treated control group. Cohort-period
estimates are aggregated to an event-study estimand and a simple weighted-average
ATT. Inference uses the multiplier bootstrap with 999 draws.

\subsection{Triple Difference}

To decompose the treatment effect into a physical-damage/displacement channel
and a belief-update channel, we estimate:

\begin{align}
  \log P_{it} &= \alpha + \beta_1 (\mathrm{Post}_t \times \mathrm{Treated}_i \times \mathrm{WUI}_i) \nonumber \\
              &\quad + \beta_2 (\mathrm{Post}_t \times \mathrm{Treated}_i)
               + \beta_3 (\mathrm{Post}_t \times \mathrm{WUI}_i) \nonumber \\
              &\quad + \beta_4 \mathrm{WUI}_i + \gamma_b + \tau_t + \varepsilon_{it}
  \label{eq:triple_diff}
\end{align}

The coefficient $\beta_1 - \beta_2$ identifies the incremental discount on
WUI-classified treated parcels relative to non-WUI treated parcels in the same
distance band and time period. Under the assumption that WUI status proxies for
forward-looking risk beliefs, this difference isolates the belief-update channel.

\subsection{Synthetic Control}

For the ZIP-level analysis, we apply three synthetic control variants:

\paragraph{ADH.} The \citet{Abadie2010} method finds convex weights $w^*$
over donor ZIP codes minimizing pre-treatment RMSPE with an importance matrix
$V$ optimized in an outer loop. Placebo inference uses the rank of the treated
unit's post-to-pre RMSPE ratio among all in-space placebos.

\paragraph{GSynth.} The generalized synthetic control \citep{Xu2017} estimates
an interactive fixed-effects counterfactual via nuclear norm minimization on the
donor panel, allowing time-varying treatment effects and heterogeneous factor
loadings.

\paragraph{AugSynth.} The augmented synthetic control \citep{BenMichael2021}
adds a ridge-regression bias correction to the ADH weights, reducing sensitivity
to imperfect pre-period fit and providing double-robustness.

Leave-one-out donor robustness is assessed for all three methods
(\texttt{src/inference/loo.py}).

%%=============================================================
\section{Spatial Spillover Analysis}
\label{sec:spatial}
%%=============================================================

\subsection{Spatial Weights}

Spatial weights $W$ are constructed as row-standardized $k$-nearest-neighbor
matrices ($k=8$) based on parcel centroid geodesic distances, using the
\texttt{SpatialWeightsFactory} in \texttt{src/spatial/weights\_phase3.py}.
All weights are stored as \texttt{scipy.sparse.csr\_matrix} objects to avoid
$O(N^2)$ memory requirements.

\subsection{Global and Local Spatial Autocorrelation}

We first test for spatial dependence in the hedonic residuals using the global
Moran's I statistic \citep{CliffOrd1981} with both the Cliff-Ord asymptotic
moment approximation and a 999-permutation empirical null distribution
(\texttt{src/esda/morans.py}).

Local clustering patterns are characterized via LISA \citep{Anselin1995},
classifying each observation as High-High (spatial hot spot), Low-Low (cold spot),
High-Low or Low-High (spatial outliers) based on the sign of the standardized
value and its spatial lag. Pseudo p-values use 999 conditional permutations
(\texttt{src/esda/lisa.py}).

\subsection{Spatial Regression Models}

We estimate three spatial regression models:

\paragraph{SAR.} The spatial autoregressive model,
$y = \rho W y + X\beta + \varepsilon$,
captures spillovers through the outcome variable. The scalar $\rho$ is
estimated by concentrated maximum likelihood with log-determinant computed
via the eigenvalue trace. Standard errors are derived from the Hessian of
the concentrated likelihood.

\paragraph{SEM.} The spatial error model,
$y = X\beta + (I - \lambda W)^{-1}\varepsilon$,
captures spatial correlation in unobservables. The scalar $\lambda$ is
estimated analogously by concentrated ML.

\paragraph{SDM.} The Spatial Durbin model,
$y = \rho W y + X\beta + WX\theta + \varepsilon$,
nests both SAR and SEM. Likelihood ratio tests evaluate the common-factor
restriction $\theta = -\rho\beta$ (which implies SAR) and the SDM-to-SEM
reduction. Model selection follows AIC and BIC (\texttt{src/spatial\_models/model\_registry.py}).

\subsection{Direct, Indirect, and Total Effects}

Following \citet{LeSage2009}, we decompose marginal effects into:
\begin{itemize}[noitemsep]
  \item \textbf{Direct}: Average own-unit response to a unit change in a covariate.
  \item \textbf{Indirect (spillover)}: Average cross-unit response to a unit change
    in a neighbor's covariate, summed over all neighbors.
  \item \textbf{Total}: Direct + Indirect.
\end{itemize}
All effects are computed via the eigenvalue trace approximation to avoid dense
matrix inversion (\texttt{src/spatial\_models/effects.py}).

\subsection{Geographically Weighted Regression}

GWR \citep{Fotheringham2002} relaxes the assumption of spatially constant
coefficients, estimating $\hat\beta(u_i,v_i)$ at each parcel location by
local weighted least squares with a bisquare adaptive kernel. The bandwidth
$h^*$ is selected by leave-one-out cross-validation using the golden-section
algorithm with pickle checkpointing (\texttt{src/gwr/bandwidth.py}).
Spatial maps of the varying treatment effect coefficient reveal geographic
heterogeneity in the price response.

%%=============================================================
\section{Results}
\label{sec:results}
%%=============================================================

\subsection{Summary Statistics}

Table~\ref{tab:summary} presents summary statistics for the estimation sample.

\begin{table}[H]
\caption{Summary Statistics}
\label{tab:summary}
\centering
\input{../docs/tables/summary_stats.tex}
\end{table}

\subsection{Hedonic Estimates}

Table~\ref{tab:hedonic} reports hedonic OLS estimates from
equation~(\ref{eq:hedonic_treatment}).

\begin{table}[H]
\caption{Hedonic Price Regression}
\label{tab:hedonic}
\centering
% \input{../docs/tables/hedonic.tex}
\textit{[TBD --- table generated by \texttt{src/outputs/tables.py}]}
\end{table}

The estimated coefficients on treatment-band indicators imply that parcels
within 0--2 km of the fire perimeter sell at a discount of approximately
[TBD]\% relative to control parcels after the fire, conditional on parcel
attributes and fixed effects.

\subsection{Callaway-Sant'Anna Event Study}

Figure~\ref{fig:event_study} plots the event-study ATT estimates
$\hat\theta^{es}(\ell)$ for event-time periods $\ell \in \{-12, \ldots, 24\}$
months relative to the August 2023 treatment date. Pre-period estimates are
statistically indistinguishable from zero (p-value = [TBD]), supporting the
parallel trends assumption. Post-treatment estimates indicate [TBD].

\begin{figure}[H]
\centering
% \includegraphics{../figures/event_study.pdf}
\textit{[TBD --- figure generated by \texttt{src/outputs/tables.py}]}
\caption{Event-Study ATT Estimates (Callaway-Sant'Anna 2021)}
\label{fig:event_study}
\end{figure}

\subsection{Triple Difference}

Table~\ref{tab:triple_diff} reports triple-difference estimates from
equation~(\ref{eq:triple_diff}).

\begin{table}[H]
\caption{Triple-Difference Decomposition}
\label{tab:triple_diff}
\centering
% \input{../docs/tables/triple_diff.tex}
\textit{[TBD --- table generated by \texttt{src/outputs/tables.py}]}
\end{table}

The estimated coefficient on the triple-interaction term ($\hat\beta_1$) is
[TBD], implying an additional discount of approximately [TBD]\% for WUI-classified
treated parcels relative to non-WUI treated parcels. The difference
$\hat\beta_1 - \hat\beta_2 = $ [TBD] is our estimate of the belief-update channel.

\subsection{Synthetic Control}

Table~\ref{tab:scm} reports the ADH, GSynth, and AugSynth treatment effect
estimates for the Maui ZIP aggregate HPI.

\begin{table}[H]
\caption{Synthetic Control Treatment Effects}
\label{tab:scm}
\centering
% \input{../docs/tables/scm_results.tex}
\textit{[TBD --- table generated by \texttt{src/outputs/scm\_tables.py}]}
\end{table}

In-space placebo inference yields a permutation p-value of [TBD] for the
ADH estimator, and leave-one-out analysis confirms the result is not driven
by any single donor unit.

\subsection{Spatial Autocorrelation}

The global Moran's I for OLS hedonic residuals is $I = $ [TBD]
($z = $ [TBD], $p < $ [TBD]), indicating significant positive spatial
autocorrelation and motivating the spatial regression analysis. LISA maps
(Figure~\ref{fig:lisa}) reveal [TBD].

\begin{figure}[H]
\centering
% \includegraphics{../figures/lisa_map.pdf}
\textit{[TBD --- figure generated by \texttt{src/outputs/spatial\_plots.py}]}
\caption{LISA Quadrant Map of Hedonic Residuals}
\label{fig:lisa}
\end{figure}

\subsection{Spatial Regression}

Table~\ref{tab:spatial} compares SAR, SEM, and SDM estimates. The preferred
model by AIC is [TBD].

\begin{table}[H]
\caption{Spatial Regression: SAR, SEM, and SDM}
\label{tab:spatial}
\centering
% \input{../docs/tables/spatial_models.tex}
\textit{[TBD --- table generated by \texttt{src/outputs/spatial\_tables.py}]}
\end{table}

Direct, indirect, and total effects from the preferred model are reported in
Table~\ref{tab:effects}. The indirect (spillover) effect of fire proximity on
parcel prices is [TBD] at the median distance from the perimeter, implying that
the aggregate market-wide price impact substantially exceeds the direct effect
on treated parcels.

\begin{table}[H]
\caption{LeSage-Pace Direct, Indirect, and Total Effects}
\label{tab:effects}
\centering
% \input{../docs/tables/effects.tex}
\textit{[TBD --- table generated by \texttt{src/outputs/spatial\_tables.py}]}
\end{table}

\subsection{GWR Spatial Heterogeneity}

Figure~\ref{fig:gwr} maps the spatially varying GWR coefficient on the treatment
indicator. The estimated discount is largest in [TBD] and smallest in [TBD],
consistent with [TBD]. The bandwidth selected by leave-one-out CV is $h^* = $ [TBD]
meters.

\begin{figure}[H]
\centering
% \includegraphics{../figures/gwr_surface.pdf}
\textit{[TBD --- figure generated by \texttt{src/outputs/spatial\_plots.py}]}
\caption{GWR: Spatially Varying Treatment Effect Coefficient}
\label{fig:gwr}
\end{figure}

%%=============================================================
\section{Conclusion}
\label{sec:conclusion}
%%=============================================================

This paper provides a multi-method assessment of the effect of the 2023 Lahaina
wildfire on residential property values in Maui County. Our hedonic and
Callaway-Sant'Anna estimates establish a causal average treatment effect of
approximately [TBD]\% on parcels within 2 km of the fire perimeter. The
triple-difference decomposition identifies a belief-update channel accounting for
approximately [TBD]\% of the total effect, representing the market's forward-looking
revision of wildfire risk for WUI-classified properties.

Synthetic control estimates at the ZIP-level HPI corroborate the parcel-level
findings and quantify the aggregate market impact on Maui real estate relative to
a synthetic counterfactual. Spatial econometric models reveal that the aggregate
price impact substantially exceeds the direct effect through spatial spillovers:
the estimated indirect spillover is [TBD] times the direct effect, implying
significant market-wide repricing of climate risk beyond the immediate destruction
zone. GWR reveals meaningful geographic heterogeneity in this response.

Taken together, these findings suggest that catastrophic wildfire events have
lasting, geographically broad effects on housing markets. The magnitude of
indirect spillovers --- the belief-update and risk-repricing components ---
exceeds the direct physical damage effect, pointing to a channel through which
climate risk propagates through insurance markets, lender risk assessments, and
buyer risk preferences. These dynamics have implications for the design of
climate adaptation policy, municipal fiscal planning in wildfire-prone areas,
and the pricing of climate risk in residential mortgage markets.

%%=============================================================
%% Appendix
%%=============================================================
\appendix

\section{Data Construction Details}

All data ingest code is in \texttt{src/ingest/}. The fire perimeter is downloaded
from the NIFC WFIGS ArcGIS FeatureServer; WUI classifications from the USFS Silvis
Lab RDS-2015-0047-3 dataset; FRED macroeconomic series via the FRED API; FHFA ZIP
HPI from the FHFA Data Tools Excel download; and Redfin neighborhood statistics
from a public S3 endpoint. A complete data source inventory with URLs and local
paths is in \texttt{docs/DATA\_SOURCES.md}.

\section{Estimation Details}

Full econometric derivations, including the concentrated ML likelihood for
SAR/SEM/SDM, the LeSage-Pace eigenvalue trace approximation for spillover effects,
the GWR golden-section bandwidth selector, and the Cliff-Ord Moran's I variance
formula, are documented in \texttt{docs/METHODOLOGY\_NOTES.md}.

%%=============================================================
%% References
%%=============================================================
\bibliography{references}

\end{document}
